\subsection{Interactive panel models}
\sourcepages{Lecture1-10.pdf, page 41, with the matrix intuition on page 40}
Consider the interactive panel model
\[
 Y_{it}=A_{i\cdot}B_{t\cdot}^{\mathsf T}+W_{it}\tau+\varepsilon_{it},
 \qquad\mathbb E[\varepsilon\mid A,B,W]=0,\tag{8.5}
\]
where $A\in\mathbb R^{n\times k}$, $B\in\mathbb R^{T\times k}$ and $k$ is a rank parameter. Equivalently,
\[
 Y_{it}=L_{it}+W_{it}\tau+\varepsilon_{it}
 \qquad\text{for a rank-$k$ matrix }L.
\]
\paragraph{My left-page matrix intuition.}
Place the $A$ values down the rows and the $B$ values across the columns. For rank one, their intersection has entry $a_i b_t$; for $k\ne1$, the entry is a sum $a_{i1}b_{t1}+a_{i2}b_{t2}+\cdots$. This reconstructs the arrows in the sketch, which lead from a row and a column to their product.

To handle this model, consider synthetic control and synthetic differences-in-differences. Suppose $W_{it}=1$ exactly when $i>n_0$ and $t>T_0$. The diagram partitions the outcome matrix at row $n_0$ and column $T_0$, with only the bottom-right block treated.

Fit the untreated units to the pretreatment trajectory of the treated group:
\[
 \sum_{i=1}^{n_0}\gamma_iY_{it}
 \approx\alpha+\frac1{n-n_0}\sum_{i=n_0+1}^nY_{it},
 \qquad t=1,\ldots,T_0.
\]
The first group is never treated; the second is treated after $T_0$. The intercept $\alpha$ plays a role similar to a fixed effect. Choose weights by
\begin{align*}
 (\widehat\gamma,\widehat\alpha)\in
 \operatorname*{arg\,min}_{\gamma',\alpha}\quad&
 \left\|\sum_{i=1}^{n_0}\gamma'_iY_{i,1:T_0}
 -\frac1{n-n_0}\sum_{i=n_0+1}^nY_{i,1:T_0}
 -\alpha\mathbf1\right\|_2^2,\\
 \text{subject to}\quad&\sum_{i=1}^{n_0}\gamma'_i=1,
 \qquad\gamma'_i\ge0.\tag{8.7}
\end{align*}
\textbf{My red annotation:} the purpose of $\gamma_i$ is to create parallel trends, so that the latent $A_i$ terms can be balanced.

Finally, estimate the treatment coefficient in the weighted two-way model:
\[
 (\widehat\tau,\widehat\alpha,\widehat\beta)
 \in\operatorname*{arg\,min}_{\tau',\alpha',\beta'}
 \sum_{i,t}\gamma_i
 (Y_{it}-\alpha'_i-\beta'_t-W_{it}\tau')^2,\tag{8.8}
\]
where $\gamma_i=1/(n-n_0)$ for $i>n_0$. The closed-form treatment coefficient is
\begin{align*}
 \widehat\tau={}&\frac1{n-n_0}\sum_{i=n_0+1}^n
 \left\{\frac1{T-T_0}\sum_{t=T_0+1}^TY_{it}
 -\frac1{T_0}\sum_{t=1}^{T_0}Y_{it}\right\}\\
 &-\sum_{i=1}^{n_0}\gamma_i
 \left\{\frac1{T-T_0}\sum_{t=T_0+1}^TY_{it}
 -\frac1{T_0}\sum_{t=1}^{T_0}Y_{it}\right\}.
\end{align*}
The red margin says this depends on $T_0$ and $n$. Under an appropriate large-panel asymptotic regime, $n_0,n_1,T\to\infty$, the estimator is consistent. A further handwritten qualification below this statement is clipped by the bottom edge of the scan; the visible fragments refer to the specification in (8.5) and asymptotic normal inference, but a complete sentence cannot be recovered.

\paragraph{Left-page postscript.}
Synthetic control is not the only method for interactive fixed effects. For example, apply $k$-means clustering to the rows of the outcome matrix $Y$, then fit a time-varying baseline model separately for each group. A short reference is written as ``p66, literature'' [page number/reference unclear].

\subsection{Identification through exchangeability}
\sourcepages{Lecture1-10.pdf, page 43, with the extension on page 42}
Recall $Y_{it}=Y_{it}(0)+W_{it}\tau$. A design-based assumption is
\[
 Y_{i\cdot}(0)\perp W_{i\cdot}\mid S_i,
 \qquad S_i=\sum_{t=1}^TW_{it}.\tag{8.10}
\]
One implication: consider a linear estimator
\[
 \widehat\tau=\sum_{i,t}\gamma_{it}Y_{it},\tag{8.11}
\]
whose weight matrix depends only on $W$. Under (8.10), the untreated potential outcomes are exchangeable within a value of $S_i$, and
\[
 \mathbb E[\widehat\tau\mid W,\gamma]
 =\sum_{i,t}\gamma_{it}\mathbb E[Y_{it}(0)\mid S_i]
 +\tau\sum_{i,t}\gamma_{it}W_{it}.
\]
Thus unbiasedness follows if
\[
 \sum_{i,t}\gamma_{it}W_{it}=1,
 \qquad
 \sum_{i:S_i=s}\gamma_{it}=0
 \quad\text{for every }t\text{ and }s\in\mathcal S,
\]
where $\mathcal S$ is the support of $S_i$. \textbf{My annotation:} within a stratum the conditional expectation is constant, so the zero sum of weights removes that term.

Choose $\gamma$ by minimizing variance under these unbiasedness constraints:
\begin{align*}
 \gamma\in\operatorname*{arg\,min}_{\gamma'}\quad&\sum_{i,t}(\gamma'_{it})^2,\\
 \text{subject to}\quad&\sum_{i,t}\gamma'_{it}W_{it}=1,
 \quad\sum_{i:S_i=s}\gamma'_{it}=0\quad\text{for all }s,t.\tag{8.14}
\end{align*}
\textbf{Blue note:} roughly, we need variation in $W_{it}$ among units with the same $S_i$.

\paragraph{Left-page extension: two sets of identifying restrictions.}
Recall the two-way model
\[
 Y_{it}=\alpha_i+\beta_t+W_{it}\tau+\varepsilon_{it},
 \qquad\mathbb E[\varepsilon\mid\alpha,\beta,W]=0.\tag{8.2, recalled}
\]
The estimator (8.11) is also unbiased under this model if
\[
 \sum_t\gamma_{it}=0\ \text{for every }i,
 \qquad\sum_i\gamma_{it}=0\ \text{for every }t,
 \qquad\sum_{i,t}\gamma_{it}W_{it}=1.
\]
\textbf{Red note:} the first two equalities cancel the fixed effects $\alpha_i$ and $\beta_t$, and the third normalizes the treatment coefficient.

Combine the restrictions to obtain a doubly robust estimator: minimize $\sum_{i,t}\gamma_{it}^2$ subject to
\[
 \sum_{i,t}\gamma_{it}W_{it}=1,
 \qquad\sum_{i:S_i=s}\gamma_{it}=0\quad\text{for all }s,t,
\]
\[
 \sum_t\gamma_{it}=0\quad\text{for all }i,
 \qquad\sum_i\gamma_{it}=0\quad\text{for all }t.
\]
It is unbiased if either (8.2) or (8.10) holds. An abbreviated reference near the top of the supplement reads ``p67, [2019]'' [the citation is not fully identifiable from the scan].

\section{Lecture 9: Instrumental variables regression}
\sourcepages{Lecture1-10.pdf, page 45, with the three IV conditions on page 44}
Instrumental variables regression addresses endogenous treatment effects, where $Y_i(0)$ and $W_i$ are dependent.

\subsection{A structural model}
Start with a constant-effect model:
\[
 Y_i(w)=Y_i(0)+w\tau,\qquad Y_i=Y_i(W_i).\tag{9.1}
\]
Unconfoundedness ordinarily means $\{Y_i(w)\}\perp W_i$; in (9.1), this is $Y_i(0)\perp W_i$. Rewrite the model as
\[
 Y_i=\alpha+W_i\tau+\varepsilon_i,\tag{9.2}
\]
where $\alpha=\mathbb E[Y_i(0)]$ and $\varepsilon_i=Y_i(0)-\mathbb E[Y_i(0)]$. Under unconfoundedness, $\mathbb E[\varepsilon_i\mid W_i]=0$, so OLS of $Y_i$ on $W_i$ estimates $\tau$.

With confounding, equation (9.2) still holds, but $\mathbb E[\varepsilon_i\mid W_i]\ne0$. OLS then targets
\begin{align*}
 \tau_{\mathrm{OLS}}
 &=\frac{\operatorname{Cov}(Y_i,W_i)}{\operatorname{Var}(W_i)}\\
 &=\frac{\operatorname{Cov}(\tau W_i+\varepsilon_i,W_i)}{\operatorname{Var}(W_i)}
 =\tau+\frac{\operatorname{Cov}(\varepsilon_i,W_i)}{\operatorname{Var}(W_i)}
 \ne\tau.
\end{align*}
The ratio is underlined as the OLS coefficient.

\subsection{Identification with an instrument}
Introduce $Z_i$:
\[
 Y_i=\alpha+W_i\tau+\varepsilon_i,\qquad\varepsilon_i\perp Z_i,
 \qquad W_i=Z_i\gamma+\eta_i.\tag{9.4}
\]
Since $Z_i$ is uncorrelated with $\varepsilon_i$,
\[
 \operatorname{Cov}(Y_i,Z_i)
 =\operatorname{Cov}(\tau W_i+\varepsilon_i,Z_i)
 =\tau\operatorname{Cov}(W_i,Z_i).
\]
Hence
\[
 \tau=\frac{\operatorname{Cov}(Y_i,Z_i)}{\operatorname{Cov}(W_i,Z_i)},
 \qquad
 \widehat\tau=\frac{\widehat{\operatorname{Cov}}(Y_i,Z_i)}
 {\widehat{\operatorname{Cov}}(W_i,Z_i)}.
\]

\paragraph{My left-page checklist for an IV.}
\begin{enumerate}
 \item \textbf{Exogenous:} $\varepsilon_i\perp Z_i$.
 \item \textbf{Relevant:} $\operatorname{Cov}(W_i,Z_i)\ne0$.
 \item \textbf{Exclusion restriction:} every effect of $Z_i$ on $Y_i$ must pass through $W_i$.
\end{enumerate}
The sketch is $Z_i\to W_i\to Y_i$; additional routes from treatment to outcome are sketched, while the instrument's influence is required to go through treatment.

\subsection{Optimal instruments}
\sourcepages{Lecture1-10.pdf, page 47, with extensions on page 46 and the estimating-equation reminder on page 48}
Allow a high-dimensional instrument $Z_i\in\mathcal Z$:
\[
 Y_i=\tau W_i+\varepsilon_i,\qquad\varepsilon_i\perp Z_i,
 \qquad Y_i,W_i\in\mathbb R.
\]
For a function $w:\mathcal Z\to\mathbb R$,
\[
 \tau=\frac{\operatorname{Cov}(Y_i,w(Z_i))}
 {\operatorname{Cov}(W_i,w(Z_i))},\tag{9.8}
\]
with nonzero denominator. Which transformation $w(\cdot)$ is optimal? The estimator is
\begin{align*}
 \widehat\tau_w
 &=\frac{\widehat{\operatorname{Cov}}(Y_i,w(Z_i))}
 {\widehat{\operatorname{Cov}}(W_i,w(Z_i))}\\
 &=\frac{n^{-1}\sum_i(Y_i-\overline Y)(w(Z_i)-\overline{w(Z)})}
 {n^{-1}\sum_i(W_i-\overline W)(w(Z_i)-\overline{w(Z)})},\tag{9.9}
\end{align*}
where $\overline Y=n^{-1}\sum_iY_i$, with analogous definitions for the other averages. It solves the estimating equation
\[
 \frac1n\sum_i(w(Z_i)-\overline{w(Z)})
 \{Y_i-\overline Y-\widehat\tau_w(W_i-\overline W)\}=0.\tag{9.10}
\]
The source writes
\[
 n(\widehat\tau_w-\tau)\Rightarrow\mathcal N(0,V_w),
 \qquad
 V_w=\frac{\operatorname{Var}(\varepsilon_i)\operatorname{Var}(w(Z_i))}
 {\operatorname{Cov}(W_i,w(Z_i))^2}.\tag{9.11, source scaling}
\]
\textit{Translator's note:} the multiplier is handwritten as $n$, but the estimating-equation reminder on page 48 uses $\sqrt n$. The usual nondegenerate normal scaling for this expression is $\sqrt n$; the mismatch is recorded explicitly.

Because $Z_i\perp\varepsilon_i$ and $\mathbb E[\varepsilon_i]=0$,
\[
 \operatorname{Var}\{\varepsilon_i(w(Z_i)-\mathbb E[w(Z_i)])\}
 =\operatorname{Var}(\varepsilon_i)\operatorname{Var}(w(Z_i)).
\]
The optimization written in the notebook is
\[
 w^*(\cdot)\in\operatorname*{arg\,min}_{w'}
 \frac{\operatorname{Cov}(W_i,w'(Z_i))^2}{\operatorname{Var}(w'(Z_i))}.
\]
Its stated solution is
\[
 w^*(z)\propto\mathbb E[W_i\mid Z_i=z].
\]
\textbf{Red interpretation:} $w^*(z)$ is the best prediction of $W_i$ from $Z_i$.
\textit{Translator's note:} minimizing $V_w$ requires \emph{maximizing}, not minimizing, the displayed covariance-squared/variance ratio. The source's ``arg min'' and its intended optimal prediction are both retained, with this discrepancy made explicit.

\paragraph{Left-page extension: conditioning on covariates.}
The condition can be extended to $\varepsilon_i\perp Z_i\mid X_i$, giving
\[
 \tau(x)=\frac{\operatorname{Cov}(Y_i,w(Z_i)\mid X_i=x)}
 {\operatorname{Cov}(W_i,w(Z_i)\mid X_i=x)}.
\]
The annotation mentions a doubly robust estimator of the average effect $\tau=\mathbb E[\tau(X)]$, and a random-forest estimator of $\tau(\cdot)$.

Another reminder on the same page: for independent $X_1,X_2$,
\begin{align*}
 \operatorname{Var}(X_1X_2)
 ={}&\operatorname{Var}(X_1)\operatorname{Var}(X_2)\\
 &+\mathbb E[X_2]^2\operatorname{Var}(X_1)
 +\mathbb E[X_1]^2\operatorname{Var}(X_2).
\end{align*}

\paragraph{Starred supplement: estimating equations and the delta method.}
If functions $\psi_i$ satisfy $\mathbb E[\psi_i(\theta)]=0$, and $\widehat\theta$ is estimated by
\[
 \frac1n\sum_{i=1}^n\psi_i(\widehat\theta)=0,
\]
then, in the regular setting described by the note's nonparametric delta-method reminder,
\[
 \sqrt n(\widehat\theta-\theta)\Rightarrow\mathcal N(0,V),
 \qquad
 V=\frac{\operatorname{Var}(\psi_i(\theta))}
 {\mathbb E[\psi_i'(\theta)]^2}.
\]

\subsection{Cross-fitting and feasible estimation}
\sourcepages{Lecture1-10.pdf, page 49}
The optimal instrument is $w^*(z)=\mathbb E[W_i\mid Z_i=z]$. The general procedure is:
\begin{enumerate}
 \item Fit a nonparametric first-stage regression to estimate $\mathbb E[W_i\mid Z_i=z]$, obtaining $\widehat w(\cdot)$.
 \item Use $\widehat w(\cdot)$ as the instrument in (9.9), obtaining $\widehat\tau_w$.
\end{enumerate}
However, with many instruments this can create overfitting bias. If $\widehat w(Z_i)$ is fit using the entire data set, it need not be independent of $\varepsilon_i$: the fitted instrument can inherit the dependence of $W_i$ on $\varepsilon_i$.

\paragraph{My blue example in the margin.}
Suppose $Z_i$ is noise. The note writes
\[
 W_i=\alpha+\beta Z_i+\eta_i,
 \qquad Y_i=\gamma+\delta\widehat W_i+\varepsilon_i
 =\gamma+\delta(\widehat\alpha+\widehat\beta Z_i)+\varepsilon_i,
\]
marks the fitted term as biased, and draws an arrow comparing it to $Y_i=\gamma+\delta W_i+\varepsilon_i$. The point of the annotation is that using an overfit first stage does not remove the problematic dependence. \textit{Translator's note:} the margin labels the two expressions ``equivalent,'' but does not supply conditions for that equivalence; it is retained as the author's warning/example rather than as a proved identity.

Use $K$-fold cross-fitting instead. With $k(i)$ the fold containing observation $i$,
\[
 \widehat\tau=
 \frac{\widehat{\operatorname{Cov}}(Y_i,\widehat w^{(-k(i))}(Z_i))}
 {\widehat{\operatorname{Cov}}(W_i,\widehat w^{(-k(i))}(Z_i))}.\tag{9.13}
\]
The corresponding fitted instrument is independent of the observation's error under the stated sampling setup. In particular, if $\widehat w^{(-k(i))}(z)$ is mean-square consistent for $\mathbb E[W_i\mid Z_i=z]$, then (9.13) is asymptotically equivalent to the oracle version of (9.9). The red reminder writes this consistency as
\[
 \lim_{n\to\infty}\mathbb E[(\widehat w_n-\mathbb E[W\mid Z])^2]=0,
\]
where $n$ is the sample size.

\subsection{Nonparametric instrumental variables regression}
\sourcepages{Lecture1-10.pdf, page 50}
The previous model assumed $Y_i=W_i\tau+\varepsilon_i$. Generalize it to
\[
 Y_i=g(W_i)+\varepsilon_i,\qquad Z_i\perp\varepsilon_i,
 \qquad Y_i,W_i\in\mathbb R,\quad Z_i\in\mathcal Z,
\]
where $g$ is a general smooth function. \textbf{Blue warning:} we cannot simply run a nonparametric regression of $Y_i$ on $W_i$, because
\[
 g(w)\ne\mathbb E[Y_i\mid W_i=w].
\]
\textit{Translator's note:} the introductory parenthesis in the scan says ``$Z_i$ acts linearly on $Y_i$,'' although the displayed earlier model is linear in $W_i$. The formula is the basis of the translation, and that wording discrepancy is recorded here.

Since $Z_i\perp\varepsilon_i$ and $\mathbb E[\varepsilon_i]=0$,
\begin{align*}
 \mathbb E[Y_i\mid Z_i=z]
 &=\mathbb E[g(W_i)+\varepsilon_i\mid Z_i=z]\\
 &=\mathbb E[g(W_i)\mid Z_i=z]
 =\int_{\mathbb R}g(w)f(w\mid z)\,dw.\tag{9.15}
\end{align*}
Here $f(w\mid z)$ is the conditional density of $W_i$ given $Z_i=z$.

A learning procedure:
\begin{enumerate}
 \item Estimate $\widehat f(w\mid z)$ using cross-fitting.
 \item Choose a suitable function class $\mathcal G$ and minimize the loss to estimate $g$:
 \[
 \widehat g\in\operatorname*{arg\,min}_{g\in\mathcal G}
 \frac1n\sum_{i=1}^n\left\{
 Y_i-\int_{\mathbb R}g(w)\widehat f^{(-k(i))}(w\mid Z_i)\,dw
 \right\}^{\!2}.\tag{9.16}
 \]
\end{enumerate}
To solve this in practice, approximate $g$ with a basis expansion,
\[
 g_K(w)=\sum_{k=1}^K\beta_k\psi_k(w),
\]
where the $\psi_k$ are predetermined basis functions. As $K$ grows, $g_K$ better approximates $g$. Then
\[
 \widehat\beta\in\operatorname*{arg\,min}_\beta
 \frac1n\sum_{i=1}^n
 \{Y_i-\widehat m^{(-k(i))}(Z_i)^{\mathsf T}\beta\}^2,\tag{9.17}
\]
with components
\[
 \widehat m_k^{(-k(i))}(Z_i)=
 \int_{\mathbb R}\psi_k(w)\widehat f^{(-k(i))}(w\mid Z_i)\,dw.
\]
\textbf{Bottom-margin interpretation:} consistency for $g(\cdot)$ needs conditions; (9.15) is complicated, while (9.17) turns the problem into regression by minimizing the loss for each basis coefficient. The blue margin adds that using (9.17) comes at the cost of regularization. The remaining bottom-edge wording is clipped and cannot be fully read. A nearby note says ``analogous to 9.13.''

\section{Lecture 10: Local average treatment effects}
\sourcepages{Lecture1-10.pdf, page 52, then the compliance table and derivation on page 51}
Recall the Lecture 9 IV model:
\[
 Y_i=\alpha+W_i\tau+\varepsilon_i,\quad\varepsilon_i\perp Z_i,
 \qquad W_i=Z_i\gamma+\eta_i.\tag{10.1}
\]
When $Z_i$ is binary,
\begin{align*}
 \tau&=\frac{\operatorname{Cov}(Y_i,Z_i)}{\operatorname{Cov}(W_i,Z_i)}\\
 &=\frac{\mathbb E[Y_i\mid Z_i=1]-\mathbb E[Y_i\mid Z_i=0]}
 {\mathbb E[W_i\mid Z_i=1]-\mathbb E[W_i\mid Z_i=0]}.\tag{10.2}
\end{align*}

\subsection{Treatment effects under noncompliance}
In an RCT with noncompliance, $Y_i\in\mathbb R$ is the ordinary outcome; $W_i\in\{0,1\}$ is the treatment actually received and is not randomized; $Z_i\in\{0,1\}$ is randomized treatment assignment.

The constant $\tau$ in (10.1) may not be appropriate, and treatment receipt may be related to the treatment effect. Define potential treatment receipt $W_i(0),W_i(1)$, with $W_i=W_i(Z_i)$ and $Y_i=Y_i(W_i(Z_i))$.

The IV validity assumptions are:
\begin{itemize}
 \item \textbf{Exclusion restriction:} $Y_i(w,z)=Y_i(w)$. Assignment affects the outcome only through treatment receipt.
 \item \textbf{Exogeneity:} $\{Y_i(0),Y_i(1),W_i(0),W_i(1)\}\perp Z_i$. Assignment is randomized.
 \item \textbf{Relevance:} $\mathbb E[W_i(1)-W_i(0)]\ne0$. Assignment affects treatment receipt.
\end{itemize}

There are four compliance types:
\begin{center}
\begin{tabular}{lcc}
 & $W_i(1)=0$ & $W_i(1)=1$\\\hline
 $W_i(0)=0$ & Never-taker & Complier\\
 $W_i(0)=1$ & Defier & Always-taker
\end{tabular}
\end{center}
Assume there are no defiers, that is, \textbf{monotonicity}:
\[
 W_i(1)\ge W_i(0)\quad\text{for every }i.
\]
Then
\begin{align*}
 &\mathbb E[Y_i\mid Z_i=1]-\mathbb E[Y_i\mid Z_i=0]\\
 &\quad=\mathbb E[Y_i(W_i(1))\mid Z_i=1]
 -\mathbb E[Y_i(W_i(0))\mid Z_i=0]\\
 &\quad=\mathbb E[Y_i(W_i(1))-Y_i(W_i(0))]\\
 &\quad=\mathbb E[\mathbf1\{C_i=\mathrm{complier}\}
 (Y_i(1)-Y_i(0))].
\end{align*}
If compliers exist, applying Bayes' rule gives
\begin{align*}
 \tau_{\mathrm{LATE}}
 &=\frac{\mathbb E[Y_i\mid Z_i=1]-\mathbb E[Y_i\mid Z_i=0]}
 {\mathbb E[W_i\mid Z_i=1]-\mathbb E[W_i\mid Z_i=0]}\\
 &=\mathbb E[Y_i(1)-Y_i(0)\mid C_i=\mathrm{complier}].\tag{10.3}
\end{align*}
This is the complier average treatment effect.

\paragraph{My blue postscript.}
When the constant-effect model (10.1) fails and there are never-takers, always-takers, defiers, and so on, the individual contrast $Y_i(1)-Y_i(0)$ is not generally observable. In particular, a never-taker is never observed with $W_i=1$. The marginal line writes an ``IATE'' expression as $\mathbb E[Y_i(1)-Y_i(0)]$ [the label and intended distinction between individual and average effects are not resolved in the scan].

\subsection{Supply and demand}
\sourcepages{Lecture1-10.pdf, page 53}
For each market $i=1,\ldots,n$, let $S_i(p,z)$ be its supply curve and $Q_i(p,z)$ its demand curve, where price $p\in\mathbb R$ and the instrument $z\in\{0,1\}$. Assume supply is continuous and increasing in price, while demand is continuous and decreasing. The equilibrium price $P_i$ is the unique solution to
\[
 S_i(P_i,Z_i)=Q_i(P_i,Z_i).
\]
Thus $Z_i$ is the IV, $P_i$ is treatment, and $Q_i(P_i,Z_i)$ is the outcome $Y_i$. The red causal chain is $Z_i\to S_i\to P_i\to Q_i(P_i)$.

Conditions:
\begin{itemize}
 \item Exclusion: $Q_i(p,z)=Q_i(p)$ for every $p,z$.
 \item Exogeneity: $\{Q_i(p),S_i(p,z)\}\perp Z_i$.
 \item Relevance: $\operatorname{Cov}(P_i,Z_i)\ne0$.
 \item Monotonicity: $S_i(p,1)\le S_i(p,0)$ for all relevant $p$ [the source writes the equilibrium-price symbol in this condition].
\end{itemize}
If $Q_i(p)$ is differentiable, let $Q_i'(p)$ denote its derivative. The author's parenthetical qualification says that ordinary differentiability is not necessary: because $Q_i(\cdot)$ is monotone decreasing, a derivative in the distributional sense exists.
\begin{align*}
 \tau_{\mathrm{LATE}}
 &=\frac{\mathbb E[Q_i\mid Z_i=1]-\mathbb E[Q_i\mid Z_i=0]}
 {\mathbb E[P_i\mid Z_i=1]-\mathbb E[P_i\mid Z_i=0]}\\
 &=\frac{\displaystyle\int
 \mathbb E[Q_i'(p)\mid P_i(0)\le p\le P_i(1)]
 P[P_i(0)\le p\le P_i(1)]\,dp}
 {\displaystyle\int P[P_i(0)\le p\le P_i(1)]\,dp}.
\end{align*}
The denominator is annotated ``total integral.''

\paragraph{Derivation in the notebook.}
\begin{align*}
 &\mathbb E[Q_i\mid Z_i=1]-\mathbb E[Q_i\mid Z_i=0]\\
 &\quad=\mathbb E[Q_i(P_i(1))\mid Z_i=1]
 -\mathbb E[Q_i(P_i(0))\mid Z_i=0]\\
 &\quad=\mathbb E[Q_i(P_i(1))-Q_i(P_i(0))]\\
 &\quad=\mathbb E\left[\int_{P_i(0)}^{P_i(1)}Q_i'(p)\,dp\right]\\
 &\quad=\int\mathbb E[Q_i'(p)\mid P_i(0)\le p\le P_i(1)]
 P[P_i(0)\le p\le P_i(1)]\,dp.
\end{align*}
\textbf{Red note:} interchange the integration order using Fubini's theorem, which is linked in the margin to exchanging the integral and expectation.

\subsection{Threshold crossing and willingness to pay}
\sourcepages{Lecture1-10.pdf, page 55, with the diagram and decomposition on page 54}
Suppose each subject has an unobserved endogenous variable $U_i$ such that
\[
 W_i=\mathbf1\{U_i\ge c(Z_i)\},
\]
where $c(z)$ is a cutoff function depending on the instrument. Without loss of generality, take $U_i\sim\operatorname{Unif}[0,1]$, giving
\[
 c(z)=1-P(W_i=1\mid Z_i=z).
\]
The IV setup imposes:
\begin{enumerate}
 \item Potential outcomes $Y_i(0),Y_i(1)$ with $Y_i=Y_i(W_i)$: the instrument has no direct route to the outcome.
 \item $\{Y_i(0),Y_i(1),U_i\}\perp Z_i$: exogeneity.
 \item $c(Z_i)$ has a nontrivial distribution: relevance.
 \item $c(z)$ is c\`adl\`ag (right-continuous with left limits) and nondecreasing.
\end{enumerate}
\textit{Translator's note:} beside the fourth condition the source also says that $W_i$ is nondecreasing in $Z_i$. With the written rule $W_i=\mathbf1\{U_i\ge c(Z_i)\}$ and increasing $c$, $W_i$ instead decreases. The source's cutoff convention is retained, and this sign inconsistency is explicitly flagged.

Define the marginal treatment effect (MTE)
\[
 \tau(u)=\mathbb E[Y_i(1)-Y_i(0)\mid U_i=u].\tag{10.6}
\]
The IV estimand is a weighted average of $\tau(u)$. In the notes, $Z_i\sim\mathcal N(0,1)$, with density $\varphi$. Let $S\subset\mathbb R$ be the discontinuity set of $c$, and
\[
 c_-(z)=\lim_{a\uparrow z}c(a).
\]
The displayed representation is
\[
 \tau_{\mathrm{LATE}}=
 \frac{\displaystyle\sum_{z\in S}
 \left\{\int_{c_-(z)}^{c(z)}\tau(u)\,du\right\}\varphi(z)
 +\int_{\mathbb R\setminus S}\tau(c(z))c'(z)\varphi(z)\,dz}
 {\displaystyle\sum_{z\in S}(c(z)-c_-(z))\varphi(z)
 +\int_{\mathbb R\setminus S}c'(z)\varphi(z)\,dz}.
\]
This displays the IV estimand as a convex average of marginal treatment effects.

\paragraph{My left-page explanation of the weights.}
The drawing plots $c(z)$ against $z$ with an open dot at the left limit and a filled dot at the value after a jump. The jump location belongs to $S$. At a jump, integrate $\tau(u)$ over the strip from $c_-(z)$ to $c(z)$ and multiply by $\varphi(z)$. On the continuous part, evaluate the treatment effect at $u=c(z)$ and weight it by $c'(z)\varphi(z)$. The arrows call $\varphi(z)$ a probability weight; the derivative represents the width contributed by the changing cutoff. The denominator adds the total weight. A second sketch shades the interval in $u$ corresponding to the jump, making the $du$ contribution explicit.

\paragraph{Bottom qualification in the source.}
The note says that a monotone $c(z)$ must have bounded variation, then writes
\[
 c(z)=c_0+\int_{-\infty}^zc'(a)\,da,
\]
where $c'$ is a nonnegative Lebesgue-integrable function.
\textit{Translator's note:} that integral-only representation describes an absolutely continuous part. It does not by itself include the jumps explicitly allowed above (nor a possible singular continuous part). The original statement is retained as a qualification requiring care, rather than silently deleting the jump terms.

\clearpage
\subsection{Single jumps, multiple jumps, and continuous cutoffs}
\sourcepages{Lecture1-10.pdf, page 57, with the derivative supplement on page 56}
\paragraph{Example 1: one jump.}
For a cutoff otherwise constant with a single jump,
\[
 c(z)=c_0+\delta_1\mathbf1\{z\ge z_1\},
\]
the source lists
\[
 \begin{array}{ll}
 \text{Never-taker:}&U_i<c_0,\\
 \text{Complier:}&c_0\le U_i<c_0+\delta_1,\\
 \text{Always-taker:}&U_i\ge c_0+\delta_1.
 \end{array}
\]
\textit{Translator's note:} these are the notebook's labels. With its increasing-cutoff convention and $W_i=\mathbf1\{U_i\ge c(Z_i)\}$, the middle group switches from treated to untreated as $z$ rises; it is therefore opposite to the earlier $W_i(1)\ge W_i(0)$ complier direction. This is the same convention inconsistency noted above.

For $K$ jumps,
\[
 c(z)=c_0+\sum_{k=1}^K\delta_k\mathbf1\{z\ge z_k\},
\]
and the note gives
\[
 \tau_{\mathrm{LATE}}=
 \frac{\displaystyle\sum_{k=1}^K
 \mathbb E[\tau(U_i)\mid c_-(z_k)\le U_i<c(z_k)]
 (c(z_k)-c_-(z_k))\varphi(z_k)}
 {\displaystyle\sum_{k=1}^K(c(z_k)-c_-(z_k))\varphi(z_k)}.
\]
It is a convex combination of the average effects in the switching strata. The weights depend on the size of each stratum in $U$ space and on $\varphi(z_k)$.
\textit{Translator's note:} the numerator's first difference appears as $z_k-c_-(z_k)$ in the handwriting, while the denominator and the immediately preceding general expression use $c(z_k)-c_-(z_k)$. The consistent cutoff difference is displayed here, with the source discrepancy recorded.

\paragraph{Example 2: a continuous cutoff.}
If $c(z)$ has no jumps, the displayed formula becomes
\[
 \tau_{\mathrm{LATE}}=
 \frac{\displaystyle\int_{\mathbb R}\tau(c(z))c'(z)\varphi(z)\,dz}
 {\displaystyle\int_{\mathbb R}c'(z)\varphi(z)\,dz}.
\]
The covariance derivation is
\begin{align*}
 \operatorname{Cov}(Y_i,Z_i)
 &=\operatorname{Cov}(Y_i(0)+(Y_i(1)-Y_i(0))W_i,Z_i)\\
 &=\operatorname{Cov}((Y_i(1)-Y_i(0))W_i,Z_i)\\
 &=\operatorname{Cov}((Y_i(1)-Y_i(0))\mathbf1\{U_i\ge c(Z_i)\},Z_i)\\
 &=\operatorname{Cov}(\tau(U_i)\mathbf1\{U_i\ge c(Z_i)\},Z_i).
\end{align*}
Define
\[
 H(z)=\mathbb E[\tau(U_i)\mathbf1\{U_i\ge c(z)\}].
\]
For standard Gaussian $Z_i$, the note uses
\[
 \operatorname{Cov}(H(Z_i),Z_i)=\mathbb E[H'(Z_i)].
\]

\paragraph{Left-page question: distributional derivatives.}
The starred supplement asks: ``Does $H'(z)$ denote the distributional derivative of $H(\cdot)$?'' It writes, in a mixed ordinary/distributional notation,
\[
 -H'(z)=
 \begin{cases}
 \displaystyle\left\{\int_{c_-(z)}^{c(z)}\tau(u)\,du\right\}\delta_z,
 &z\in S,\\
 \tau(c(z))c'(z),&\text{otherwise}.
 \end{cases}
\]
Here $\delta_z$ is a Dirac delta at $z$. The accompanying reminder gives the informal defining properties of a delta centered at zero: $\delta(x)=0$ for $x\ne0$ and $\int_{-\infty}^{\infty}\delta(x)\,dx=1$. The question mark and the distributional nature of this supplement are retained.

\subsection{Estimating the marginal treatment effect}
\sourcepages{Lecture1-10.pdf, page 58}
At a smooth point of $c(z)$,
\[
 \tau(c(z))=
 \frac{\displaystyle\frac d{dz}\mathbb E[Y_i\mid Z_i=z]}
 {\displaystyle\frac d{dz}P(W_i=1\mid Z_i=z)}.\tag{10.12}
\]
Assume the derivatives are well behaved, the numerator and denominator can be estimated, and the denominator is nonzero.

\paragraph{Proof in the notebook.}
\begin{align*}
 \mathbb E[Y_i\mid c(Z_i)=c]
 &=\mathbb E[Y_i(0)+\mathbf1\{U_i\ge c\}(Y_i(1)-Y_i(0))
 \mid c(Z_i)=c]\\
 &=\mathbb E[Y_i(0)+\mathbf1\{U_i\ge c\}(Y_i(1)-Y_i(0))]\\
 &=\mathbb E[Y_i(0)]+\int_c^1\tau(u)\,du.
\end{align*}
The last step is marked ``Fubini.'' Therefore
\[
 \tau(c)=-\frac d{dc}\mathbb E[Y_i\mid c(Z_i)=c].
\]
By the chain rule,
\[
 \frac d{dz}\mathbb E[Y_i\mid Z_i=z]
 =\left.\frac d{dc}\mathbb E[Y_i\mid c(Z_i)=c]\right|_{c=c(z)}c'(z).
\]
Since $U_i\sim\operatorname{Unif}(0,1)$ independently of $Z_i$,
\[
 c'(z)=-\frac d{dz}P(W_i=1\mid Z_i=z).
\]
Combining these identities gives (10.12), as checked in the original.
